(1) 在 $(0,0,-h)$ 放置虚构镜像电荷 $-q$，则对 $z>0$
\[
\phi(\rho,z)=\frac{q}{4\pi\epsilon_0}\left[
\frac1{\sqrt{\rho^2+(z-h)^2}}-
\frac1{\sqrt{\rho^2+(z+h)^2}}\right].
\]
在 $z=0$ 两项严格抵消，符合接地边界。实际电荷位置处仍保留 Poisson 方程源项，在其余 $z>0$ 区域满足 Laplace 方程；由静电唯一性定理该势是所求物理解。

(2) 在导体外侧 $z=0^+$，法向外电场为 $E_z=-\partial_z\phi$，表面自由电荷密度为 $\sigma=\epsilon_0E_z$。计算得到
\[
\sigma(\rho)=-\frac{qh}{2\pi(\rho^2+h^2)^{3/2}}<0.
\]
将圆环面元 $2\pi\rho\,d\rho$ 积分，
\[
Q_{\rm ind}=\int_0^\infty\sigma\,2\pi\rho\,d\rho
=-qh\int_0^\infty\frac{\rho\,d\rho}{(\rho^2+h^2)^{3/2}}
=-q.
\]
因此诱导面电荷总量恰为 $-q$，尽管分布延伸到无限远。

(3) 实际点电荷受导体的力等效于位于距离 $2h$ 的镜像电荷的库仑力，指向导体：
\[
\boldsymbol F=-\frac{q^2}{16\pi\epsilon_0h^2}\hat z.
\]
如果把镜像当作一颗真实独立外部电荷而直接取 $q\phi_{\rm image}$，会把能量算大两倍，因为接地导体上的诱导电荷随实际电荷位置同步重排。正确的相互作用势能是
\[
U(h)=\frac12q\phi_{\rm image}(0,0,h)
=-\frac{q^2}{16\pi\epsilon_0h},
\]
满足 $-\partial_hU=F_z$。

(4) 以无穷远相互作用能为零，缓慢移走点电荷时所需外力功是
\[
W_{\rm ext}=U(\infty)-U(h)
=\frac{q^2}{16\pi\epsilon_0h}>0.
\]
当 $h\to\infty$，力和所需功均趋零；$h\to0^+$ 时连续介质理想导体模型给出发散趋势，但当距离进入材料微观尺度时，点电荷和完美导体模型将不再可靠，不可声称物理体系真有无限可提取能量。
